\scanpage{18}
% archive foliation: 17
On désigne par $\overline P_i$ \inserted{$\subset Q_i$}
l'image inverse \inserted{dans $Q_i$} du
sous-groupe de torsion de $Q_i/P_i\simeq H_{i-1}$,
$N_{i-1}=Q_i/\overline P_i$ est de t.f. sans torsion, donc
projectif de t.f., donc $Q_i\simeq\unsure{\overline P_i|P_i}\oplus N_{i-1}$
\margin{$\simeq H_{i-1}/\mathrm{torsion}$}
et le complexe $K_i$ est isom. à la
somme directe des complexes
\diagram{\begin{tikzcd}[column sep=large]
\text{deg} &&&\text{deg $i$}&i-1&&\\
\overline K_i:\quad &\cdots\arrow[r,dash]&0\arrow[r]&P_i\arrow[r,hook]&\overline P_i\arrow[r]&0\arrow[r]&0\\
N_i[-i]\quad &&0\arrow[r]&0\arrow[r,dash]&N_{i-1}\arrow[r]&0&
\end{tikzcd}}

Enfin, $P_i\subset\overline P_i$ est une inclusion de
modules projectifs de t.f. tels que
$\overline P_i/P_i$ $(\simeq\mathrm{tors}\ H_{i-1})$ soit de torsion,
et est justiciable de la théorie des
diviseurs élémentaires: \inserted{si $k$ \uline{principal}}
\deleted{On aura}
ce complexe est une somme directe
\deleted{\unsure{$\overline K_i\simeq$}}
de complexes de la forme
\diagram{\begin{tikzcd}
0\arrow[r,dash]&\underset{i}{k}\arrow[r,"n_{i\alpha}"]&
\underset{i-1}{k}\arrow[r]&0
\end{tikzcd}}
\margin{$1\leq\alpha\leq N_i$
$=\operatorname{rg}P_i=\operatorname{rg}\overline P_i$}
avec les $n_{i\alpha}\in k-\{0\}$ (les $n_{i\alpha}$ pouvant être
une unité --- donc $1$ alors); le
nombre des $n_{i\alpha}\in k^*$ qui interviennent
est noté $\mu_i$ (il leur \unsure{correspond|corresponde}
les facteurs $(\Lambda^i\Phi)^{\mu_i}$), et si on
prend les \inserted{autres} $n_{i,\alpha}$ tels qu'ils sont croissants
pour la divisibilité, on trouve une suite
unique; les multiplicités \inserted{d'occurrence} des diviseurs $\ell$
